\author{OTS}
\documentclass[14pt]{../../inc/mybib}

\setincpath{../../inc/}

\usepackage{bible_style}
\usepackage{textmakro}
\graphicspath{{../../assets/images/}}
\usepackage{header}

\newcommand{\Name}{Fredy}

% ensure scrlayer-scrpage has sufficient footheight
\setlength{\footheight}{10.4pt}

\begin{document}
\begin{bibelbox}{SCHL}{Kol}{3:15-17}
    Lasst das Wort des Christus reichlich in euch wohnen in aller Weisheit; lehrt und ermahnt einander und singt mit Psalmen und Lobgesängen und geistlichen Liedern dem Herrn lieblich in eurem Herzen. Und was immer ihr tut in Wort und Werk, das tut alles im Namen des Herrn Jesus und dankt Gott, dem Vater, durch ihn.
\end{bibelbox}
\section{Begrüssung}
Ich möchte auch euch alle recht herzlich zum Gottesdienst begrüssen. Schön, dass ihr gekommen seid, um unserem \herr N zu danken, ihn zu loben und zu preisen.

Auch dich \Name{} will ich herzlich begrüssen. Schön, dass du wieder bei uns bist und mit uns Gottesdienst feierst.
\begin{itemize}
    \item[] \beten{}
    \item[] Und anschliessend singen wir zusammen das Lied
    \item[] \lied{Heilig, heilig, heilig, Gott, Du bist heilig}
\end{itemize}

\section{Informationen}
\begin{itemize}
    \item \bt[infos]{Bibel und Gebetsabend:} Do, 23.April.2026 20:00 Uhr Bibel und Gebetsabend mit  Markus 6,45-
    \item \bt[infos]{Nächster Gottesdienst:} So, 26.April.2026 14:45 Uhr mit Qumran und Jesaja        
    \item \bt[infos]{Kollekte:} Die Kollekte wird hinten in der grünen Box gesammelt und wird vollumfänglich für den Bau dieser Gemeinde eingesetzt.
\end{itemize}

\section{ Input }
\begin{spacing}{1.5}
    \begin{block}[Gott hat die Kontrolle! Aber wer ist Gott?]
        Ein halber Jahr lang habe ich mit euch über den Heilsplan Gottes gesprochen. Viele male habe ich gesagt, dass die Kontrolle hat, das Gott suverän ist und dass er alles leitet.

        Habt ihr euch nicht ab und zu gefragt, wer denn eigentlich dieser Gott ist? Es ist ja so schön einfach zu sagen, dass Gott die Kontrolle hat. Oft merken wir aber nichts davon.
    \end{block}
    \begin{block}[Ende]
        So das war der letzte Teil unserer Heilsgeschichte Gottes. Ich habe euch hier ein Ablauf der Heilsgeschichte Gottes ausgedruckt. Wir mussten diesen in der Bibelschule lernen und es ist eine schöne Übersicht, über das was ich euch im letzten halben Jahr erzählt habe. Nächsten Sonntag fangen ich ein neues Thema an, weiss aber noch nicht genau was. Wenn jemand etwas auf dem Herzen hat, dann könnt ihr mir das gerne sagen. 
    \end{block}
    \lied{Allein deine Gnade genügt}.
\end{spacing}
\section{Predigt}
\textbf{Nach der Predigt}\newline
Danken für die Predigt.
% \section{Abendmahl}
% \begin{itemize}
%     \item [] Beten für das Brot
%     \item [] Beten für den Wein
% \end{itemize}
\section{Abschluss}
\begin{itemize}
    \item [] \lied{Wenn Friede mit Gott}.
    \item [] \beten{}
\end{itemize}
\begin{bibelbox}{SCHL}{IIKor}{13:13}
    Die Gnade des Herrn Jesus Christus und die Liebe Gottes und die Gemeinschaft des Heiligen Geistes sei mit euch allen! Amen
\end{bibelbox}
\end{document}